\documentclass{beamer}
\usepackage{graphicx,color,amsmath,amssymb,amsfonts,amsbsy,amsthm}
% Node the following tag will orient things right using dvi2pdf.
% Process with LaTeX first, then dvi2pdf.
\special{landscape}
% Fuzz -----------------------------------------------------------
\hfuzz5pt % Don't bother to report overfull boxes < 5pt
% MATH -----------------------------------------------------------
\input{kalman_cmdmaster.tex}
% THEOREM ENV. ---------------------------------------------------
\newtheorem{thm}{Theorem.}
\newtheorem{cor}[thm]{Corollary.}
\newtheorem{con}[thm]{Conjecture.}
\newtheorem{lem}[thm]{Lemma.}
\newtheorem{prop}[thm]{Proposition.}
\newtheorem{defn}[thm]{Definition.}
\renewcommand{\thethm}{}  % unnumbered
\pagestyle{plain}
%%% --------------------------------------------------------------
\begin{document}
\title{
\vfill \hrulefill \vfill
\begin{center}
Overview of Kalman Filter Theory and Navigation Applications\\
\end{center}
}
\author{Michael L. Carroll}
\date{Feb 27, 2004\\
\vfill \hrulefill \vspace*{.25in} \vfill
\begin{tiny}
\copyright 2003 by Michael L. Carroll
\end{tiny}
} \maketitle
% ----------------------------------------------------------------
\begin{frame}
\textbf{Day 5: Details of the Combat Talon I Navigation Kalman
Filter}
 \begin{itemize}
 \item Segment 1:  The Combat Talon I States
 \item Segment 2:  The Combat Talon I Navaids
 \item Segment 3:  The Combat Talon I Kalman Filter High-Level Software Architecture
 \item Segment 4:  The Combat Talon I Kalman Filter JOVIAL Implementation
\end{itemize}
\end{frame}
\input{day5_seg1_ct_states/ct_states_main.tex}
\input{day5_seg2_ct_navaids/ct_navaids_main.tex}
% ----------------------------------------------------------------
\end{document}
